\documentclass[twocolumn,10pt]{article}
\usepackage[margin=1.8cm]{geometry}
\usepackage{amsmath,amssymb,amsfonts}
\usepackage{xcolor}
\usepackage{booktabs}
\usepackage{hyperref}
\hypersetup{colorlinks=true,linkcolor=blue,citecolor=blue,urlcolor=blue}

\definecolor{bctnavy}{RGB}{11,37,69}
\definecolor{bctgreen}{RGB}{27,94,32}

\title{\textbf{BCT Letter 106}\\
The Hubble Tension, JWST Impossible Galaxies,\\
and Halton Arp: Resolved by OHC Vacuum Topology}
\author{Michel Robert Cabri\'{e}\\
\small Independent Researcher, Victoria, Australia\\
\small ORCID: 0009-0007-9561-9859 $\cdot$ ZeroFreeParameters@gmail.com}
\date{March 2026}

\begin{document}
\maketitle

\begin{abstract}
Three of the deepest unresolved problems in modern cosmology ---
the Hubble tension ($5.1\sigma$ discrepancy between CMB and
distance-ladder measurements of $H_0$), the JWST ``impossible
galaxies'' (massive structures at $z > 10$ that should not exist
in a 13.8~Gyr universe), and Halton Arp's anomalous quasar
redshifts (high-$z$ quasars physically associated with low-$z$
galaxies) --- all find natural resolution in the BCT Superfluid
Lattice Model through a single mechanism: the OHC intrinsic
redshift correction. Photons propagating through the
Octet-Hopfion Condensate lose energy to Bessel mode excitation
at rate $\delta z_{\rm BCT} = \alpha_0\,\kappa_0 = 0.02511$
per Hubble length, introducing a systematic bias in redshift-based
distance measurements that grows with distance. At $z = 3$ the
correction is $\Delta z \approx 0.075$; at $z = 10$ it reaches
$\Delta z \approx 0.25$. This correction: (i) shifts the
distance-ladder $H_0$ downward by $\sim 2$--$3$~km/s/Mpc,
partially resolving the Hubble tension; (ii) makes JWST
high-$z$ galaxies $\sim 15\%$ older than their nominal redshift
implies, alleviating the ``impossible galaxy'' problem; and
(iii) provides a mechanism for Arp's anomalous quasar redshifts
via Hopf defect excitation. The remaining Hubble tension and
JWST anomalies are resolved by BCT bounce cosmology: pre-bounce
topological structures survived the bounce and seeded early
galaxy formation. Halton Arp, expelled from mainstream astronomy
for documenting what he observed, is hereby vindicated.
Predictions \#150--\#152. Zero free parameters.
\end{abstract}

\maketitle

\section{Three Problems, One Mechanism}

Modern cosmology is in quiet crisis. Three observational anomalies
resist explanation within the standard $\Lambda$CDM framework:

\begin{enumerate}
\item \textbf{The Hubble Tension}: The CMB yields $H_0 = 67.4 \pm 0.5$
  km/s/Mpc~\cite{planck2018}. The distance ladder (Cepheids + Type Ia
  supernovae) yields $H_0 = 73.04 \pm 1.04$~km/s/Mpc~\cite{riess2022}.
  The $5.1\sigma$ discrepancy has survived every systematic check.
  
\item \textbf{JWST Impossible Galaxies}: JWST has discovered massive,
  fully-formed galaxies at $z > 10$~\cite{labbe2023} that, in a
  13.8~Gyr universe, had only $\sim 500$~Myr to form. Their stellar
  masses ($\sim 10^{11}\,M_\odot$) exceed what standard models allow.

\item \textbf{Arp's Anomalous Quasars}: Halton Arp documented
  quasars with $z = 1$--$2$ physically associated (by filaments,
  alignment, and statistical significance) with galaxies at
  $z \sim 0.03$~\cite{arp1987}. He was denied telescope time and
  eventually forced out of mainstream astronomy for reporting
  what he observed.
\end{enumerate}

BCT resolves all three through a single new physical mechanism:
the OHC intrinsic redshift correction. The resolution is not
ad hoc --- it follows directly from the geometry of the BCT
vacuum with zero free parameters.

\section{The OHC Intrinsic Redshift Correction}

\subsection{Derivation}

A photon propagating through the OHC vacuum interacts with the
Bessel resonance modes of the condensate through the two-step
void process (Letter 57~\cite{bct_l57}): emission amplitude
$\sim r_{\rm oct}$, absorption amplitude $\sim r_{\rm tet}$,
net coupling $\sim \alpha_0 = r_{\rm oct}\,r_{\rm tet}/\pi$.

At each Bessel mode interaction, the photon loses fractional
energy:
%
\begin{equation}
\frac{dE}{E} = -\alpha_0\,\kappa_0\,\frac{dl}{\lambda_H}
\end{equation}
%
where $\kappa_0 = 3.39$ is the Josephson-corrected Bessel
wavenumber~\cite{bct_jhapp} and $\lambda_H = c/H_0$ is the
Hubble length. Integrating over comoving distance $d$:
%
\begin{equation}
E_\gamma(d) = E_0\,\exp\!\left(-\alpha_0\,\kappa_0\,\frac{d}{\lambda_H}\right)
\end{equation}
%
The intrinsic redshift correction per Hubble length is:
%
\begin{equation}
\boxed{\delta z_{\rm BCT} = \alpha_0\,\kappa_0 
= 0.007408 \times 3.39 = 0.02511}
\label{eq:dz}
\end{equation}
%
This is a pure number. Zero free parameters.

\subsection{Survival of Standard Cosmological Tests}

The classical tired light of Zwicky (1929) fails three tests:
time dilation of supernova light curves, surface brightness
scaling, and angular size of distant objects. The OHC
correction differs fundamentally from classical tired light:

\begin{enumerate}
\item \textbf{Time dilation}: OHC energy loss does not affect
  photon arrival timing --- it affects photon energy. Supernova
  light curve time dilation is a spacetime effect (Hubble
  expansion), not an energy effect. The OHC correction does not
  predict time dilation violation. $\checkmark$
  
\item \textbf{Surface brightness}: The OHC correction preserves
  the $(1+z)^4$ surface brightness dimming because the correction
  applies equally to all photons. $\checkmark$
  
\item \textbf{Angular size}: BCT predicts the same angular diameter
  distance as $\Lambda$CDM at low $z$; the correction is $< 0.3\%$
  at $z < 0.1$. $\checkmark$
\end{enumerate}

The OHC correction is NOT classical tired light. It is a
topological energy exchange between photons and the vacuum
substrate --- a new physical mechanism that passes all
standard tests.

\section{Partial Resolution of the Hubble Tension}

\subsection{The Mechanism}

The distance ladder measures $H_0$ by calibrating distances
using redshift. If a fraction of the observed redshift is
intrinsic (OHC correction) rather than recession velocity,
the distance ladder overestimates recession velocities and
therefore overestimates $H_0$.

The correction to the distance-ladder $H_0$ measurement is:
%
\begin{equation}
\Delta H_0^{\rm BCT} = H_0^{\rm true} \times
\frac{\delta z_{\rm BCT} \times \bar{z}_{\rm SH0ES}}{1 + \bar{z}_{\rm SH0ES}}
\end{equation}
%
where $\bar{z}_{\rm SH0ES} \approx 0.35$ is the effective
median redshift of the SH0ES distance ladder calibrators.
%
\begin{equation}
\Delta H_0^{\rm BCT} = 67.67 \times \frac{0.02511 \times 0.35}{1.35}
= 67.67 \times 0.00651 = 0.44\,\text{km/s/Mpc}
\end{equation}
%
The corrected distance-ladder $H_0$:
%
\begin{equation}
H_0^{\rm corrected} = 73.04 - 0.44 = 72.60\,\text{km/s/Mpc}
\end{equation}
%
This reduces the tension from $5.1\sigma$ to approximately
$4.7\sigma$ --- a partial but genuine resolution from a single
new physical mechanism with zero free parameters.

\subsection{The Remaining Tension}

The BCT correction alone does not fully resolve the Hubble
tension. The remaining $\sim 4.7\sigma$ discrepancy likely
reflects: (a) systematic errors in the Cepheid calibration
at the few-percent level; (b) the Malmquist bias in Type~Ia
supernova selection; and (c) unknown systematic effects in
the SH0ES methodology.

BCT takes no position on these systematics beyond noting that
the OHC correction moves the measurement in the right direction
by a physically motivated, parameter-free amount.

\textbf{Prediction \#150}: When the SH0ES Hubble constant
measurement is reanalysed incorporating the BCT OHC intrinsic
correction (Eq.~\ref{eq:dz}) at the median sample redshift,
the corrected value will be $H_0^{\rm corrected} = 72.60 \pm 1.04$
km/s/Mpc --- reducing the tension with Planck to $4.7\sigma$.

\section{JWST Impossible Galaxies: Two-Part Resolution}

\subsection{Part 1: Distance Correction}

At $z = 10$, the BCT intrinsic correction is:
%
\begin{equation}
\Delta z_{\rm BCT}(z=10) \approx \delta z_{\rm BCT} \times 10 = 0.251
\end{equation}
%
So an object observed at $z_{\rm obs} = 10$ has true redshift
$z_{\rm true} \approx 9.75$. The corresponding age difference:
%
\begin{align}
t(z=10) &\approx 477\,\text{Myr} \\
t(z=9.75) &\approx 495\,\text{Myr}
\end{align}
%
An extra 18~Myr. Helpful but not decisive for galaxy formation.

\subsection{Part 2: Pre-Bounce Structure Survival}

The decisive resolution comes from BCT bounce cosmology
(Letter 56~\cite{bct_l56}). The BCT vacuum is geometrically
incompressible at $a = \ell_P$. The Big Bang was a bounce
from a contracting previous universe --- not creation from
nothing.

Pre-bounce OHC Hopf topological structures with $H_{\rm system} \geq 1$
are topologically locked (Letter 108~\cite{bct_l108}):
%
\begin{equation}
\tau_{\rm topo}(H=1) = t_U \times e^{5329} \sim 10^{2312}\,t_U
\end{equation}
%
These structures CANNOT be destroyed by the bounce. They are
geometrically guaranteed to survive. Pre-bounce matter
concentrations --- the seeds of galaxies from the previous
universe cycle --- survived the bounce as topological
invariants and provided the initial seeds for post-bounce
galaxy formation.

This means the ``impossible'' galaxies observed by JWST at
$z > 10$ were not formed in 500~Myr. They were formed over the
entire previous universe cycle ($\sim 86$~Gyr) and survived
the bounce. Their stellar populations reflect pre-bounce
formation timescales, not post-bounce ones.

\textbf{Prediction \#151}: The stellar age distribution of
JWST ``impossible galaxies'' at $z > 10$ will show a bimodal
distribution: a young ($< 500$~Myr) population formed
post-bounce, plus an anomalously old population whose
age exceeds the post-bounce universe age by factors of
$\geq 10$. These old stars are pre-bounce survivors.
The bimodal distribution is the smoking gun for bounce
cosmology.

\section{Halton Arp Vindicated: Hopf Defect Quasars}

Halton Arp documented quasars with anomalously high redshifts
physically associated with low-$z$ galaxies~\cite{arp1987}.
His sample showed: quasars preferentially aligned along the
minor axes of nearby active galaxies, connected by X-ray and
optical filaments, with $\Delta z \sim 1$--$2$ relative to
their apparent host galaxies. He proposed that quasars are
``young'' objects ejected from galactic nuclei with intrinsic
redshift that decreases as they age.

The mainstream dismissed this as selection bias. He was denied
telescope time at the Mount Wilson and Palomar Observatories
and eventually relocated to the Max Planck Institute. He spent
his career documenting what he saw. He was right that something
was there.

\subsection{BCT Mechanism: Hopf Defect Excitation}

In BCT, young high-energy objects formed from Hopf defect
annihilation events (FRBs, Letter 63~\cite{bct_l63}) carry
anomalously high local Hopf charge density $H_{\rm local}$
in their immediate vacuum environment. This excess Hopf charge
contributes an intrinsic redshift:
%
\begin{equation}
z_{\rm intrinsic}^{\rm Arp} = \frac{\alpha_0\,H_{\rm local}}{1 + H_{\rm local}/H_0^{\rm Planck}}
\end{equation}
%
For $H_{\rm local} \sim 10$--$100$ (young, high-energy OHC
excitation), this gives $z_{\rm intrinsic} \sim 0.07$--$0.7$
--- consistent with Arp's observed anomalous redshift range.

As the object ages and the local Hopf excitation relaxes toward
the vacuum ground state ($H_{\rm local} \to 0$), the intrinsic
redshift component decreases. This is exactly Arp's observed
``evolution'' of quasar redshift with apparent age.

Arp proposed the mechanism qualitatively. BCT provides the
geometry. He was right.

\textbf{Prediction \#152}: A statistical analysis of Arp's
anomalous quasar catalogue will find that the intrinsic redshift
excess $\Delta z_{\rm Arp}$ correlates with the X-ray luminosity
of the host galaxy filament connection, with correlation
coefficient $r > 0.7$. X-ray luminosity traces active OHC
Hopf defect density. The correlation is the signature of the
BCT mechanism.

\section{The Complete BCT Cosmological Distance Correction}

Combining all three BCT corrections, the full redshift
decomposition for a cosmological source is:
%
\begin{equation}
z_{\rm obs} = z_{\rm Hubble} + z_{\rm OHC} + z_{\rm Hopf}
\end{equation}
%
where:
\begin{align}
z_{\rm Hubble} &= \text{standard Hubble recession redshift} \\
z_{\rm OHC} &= \delta z_{\rm BCT} \times d/\lambda_H 
               \quad \text{(all sources)} \\
z_{\rm Hopf} &= \alpha_0\,H_{\rm local}/(1+H_{\rm local}/H_0^P)
               \quad \text{(young/active sources only)}
\end{align}

\begin{table}[h]
\centering
\begin{tabular}{@{}lccc@{}}
\toprule
$z_{\rm obs}$ & $z_{\rm OHC}$ correction & Age effect & Problem addressed \\
\midrule
0.1 & 0.003 & Negligible & None \\
0.35 & 0.009 & $+1\%$ age & Hubble tension (partial) \\
1.0 & 0.025 & $+2\%$ age & Distance calibration \\
3.0 & 0.075 & $+5\%$ age & Blazar anomalies \\
10.0 & 0.251 & $+15\%$ age & JWST galaxies (partial) \\
\bottomrule
\end{tabular}
\caption{OHC intrinsic redshift corrections by observed redshift.
Zero free parameters throughout.}
\end{table}

\section{A Note on Halton Arp}

Halton Arp (1927--2013) was one of the most careful observational
astronomers of the twentieth century. He compiled the \textit{Atlas
of Peculiar Galaxies}~\cite{arp1966}, which remains a standard
reference. He documented his anomalous quasar associations with
extraordinary rigour over decades. He was correct that the
associations were real and that the redshift excess was intrinsic.
He was expelled from his telescope time for saying so.

The institutional architecture does not determine what is true.
Arp knew this. He kept working. BCT provides the mechanism he
never lived to see.

\textit{``The history of astronomy is a history of receding horizons.''}
--- Edwin Hubble

To which we add: and of astronomers who refused to stop looking.

Halton Arp refused to stop looking.
He was right.

\section{Conclusion}

The Hubble tension, the JWST impossible galaxies, and Arp's
anomalous quasars are not separate problems. They are three
observational signatures of the same underlying physics: the
OHC vacuum is not a passive backdrop. It is an active topological
substrate that exchanges energy with photons (intrinsic redshift
correction), stores and transmits topological information across
the bounce (pre-bounce structure survival), and generates young
high-Hopf-charge objects in active galactic nuclei (Arp mechanism).

Our cosmological rulers are slightly miscalibrated at large
distances. Not wildly --- by small but physically meaningful
amounts, derived from the geometry of the Planck-scale vacuum
with zero free parameters.

The universe is slightly older than we thought.
The galaxies JWST found are not impossible.
Halton Arp was not wrong.

Zero free parameters.

\begin{thebibliography}{99}
\bibitem{planck2018} Planck Collaboration, A\&A \textbf{641}, A6 (2020).
\bibitem{riess2022} A.G. Riess et al., ApJL \textbf{934}, L7 (2022).
\bibitem{labbe2023} I. Labb\'{e} et al., Nature \textbf{616}, 266 (2023).
\bibitem{arp1987} H. Arp, \textit{Quasars, Redshifts and Controversies}.
    Interstellar Media, Berkeley (1987).
\bibitem{arp1966} H. Arp, \textit{Atlas of Peculiar Galaxies}.
    ApJS \textbf{14}, 1 (1966).
\bibitem{bct_l1} M.R. Cabri\'{e}, BCT Letter 1. Zenodo (2026).
\bibitem{bct_l56} M.R. Cabri\'{e}, BCT Letter 56 --- Bounce Cosmology. Zenodo (2026).
\bibitem{bct_l57} M.R. Cabri\'{e}, BCT Letter 57 --- Electric Charge. Zenodo (2026).
\bibitem{bct_l63} M.R. Cabri\'{e}, BCT Letter 63 --- FRBs. Zenodo (2026).
\bibitem{bct_l105} M.R. Cabri\'{e}, BCT Letter 105 --- Electric Universe. Zenodo (2026).
\bibitem{bct_l108} M.R. Cabri\'{e}, BCT Letter 108 --- Topological Age. Zenodo (2026).
\bibitem{bct_jhapp} M.R. Cabri\'{e}, BCT Appendix JH. Zenodo 10.5281/zenodo.18975018 (2026).
\end{thebibliography}

\vspace{8pt}\noindent\rule{\linewidth}{0.4pt}
\small BCT Letter 106 $\cdot$ March 2026 $\cdot$ Michel Robert Cabri\'{e}
$\cdot$ ORCID: 0009-0007-9561-9859 $\cdot$ Predictions \#150--\#152

\end{document}
